\subsection{Censo de calendarios y compatibilidad simultánea de las biografías regionales}
Las filas anteriores se reúnen, sin aplicar todavía ningún lector decimal,
en las tres biografías de bloques producidas por las transiciones:
\begin{align}
 B^{\rm cl}
  &=010211|012222|010211|002111|110221,\notag\\
 B^{\rm pr}
  &=201101|121221|102011|012222|102011,
 \label{act21:eq:biografias-tpk-tres-caracteres}\\
 B^{\rm au}
  &=121200|112202|121020|010210|010200.\notag
\end{align}
En efecto, los pares consecutivos de los dos primeros tramos son las filas
de \(X_0,Y_0\), y los de los dos últimos son las filas de \(X_1,Y_1\).
Así, las biografías son otra escritura del registro de doce transiciones, no
una colección de prefijos convencionales.

Para comprobar la palabra operativa del calendario, sea
\(c=c_1c_2c_3c_4\in\{0,1\}^4\) y, para cada régimen
\(\chi\in\{{\rm cl},{\rm pr},{\rm au}\}\), defínase
\[
 u_0^\chi=b_0^\chi,\qquad
 u_j^\chi=u_{j-1}^\chi L_{c_j}\quad(1\le j\le4).
\]
Para condensar el censo, denotemos por
\[
 N_{\mathrm{ar}}(c):=
 \#\{(j,\chi):u_j^\chi=b_j^\chi,\ 1\le j\le4\},\qquad
 N_{\mathrm{bio}}(c):=
 \#\{\chi:(u_0^\chi,\ldots,u_4^\chi)=B^\chi\}.
\]
La multiplicación exacta en \(\FF_3\) de los dieciséis calendarios da:
\[
\begin{array}{c|c|c}
 c&N_{\mathrm{ar}}(c)&N_{\mathrm{bio}}(c)\\ \hline
0000,0001&6&0\\
0010&9&0\\
0011&12&3\\
0100,0101,0110,0111&3&0\\
1000,1001,1010,1011&0&0\\
1100,1101,1110,1111&0&0
\end{array}
\tag{D}
\label{act21:eq:censo-dieciseis-calendarios}
\]
La tabla agota \(2^4=16\) posibilidades y cada una de sus entradas se obtiene
por cuatro multiplicaciones de una fila por una de las dos matrices impresas.
Por consiguiente, \(0011\) es el único calendario que reconstruye
simultáneamente las tres biografías completas. En notación operatoria,
\begin{equation}
 (L_0,L_0,L_1,L_1)
 \label{act21:eq:calendario-0011-interno}
\end{equation}
no es una prescripción independiente: es la única lectura binaria compatible
con las doce aristas ya producidas por \(U_t\).

El calendario \eqref{act21:eq:calendario-0011-interno} produce cuatro nuevos
bloques de seis componentes. Sus concatenaciones forman
\[
 w_6\longrightarrow w_{12}\longrightarrow w_{18}
 \longrightarrow w_{24}\longrightarrow w_{30}.
\]
Los subíndices de \(w\) indican longitud acumulada. La frontera \(R_{36}\) conserva la estructura terminal y abre la relación de prolongación; no se renombra como un último bloque periódico que pudiera repetirse indefinidamente.

La identidad \(L=X^{-1}Y\) demuestra unicidad una vez fijadas las
transiciones; no sustituye su producción por \(U_t\). Recíprocamente,
\eqref{act21:eq:censo-dieciseis-calendarios} prueba dentro del artículo la unicidad
del orden de los cuatro transportes. La cadena causal queda así cerrada:
el censo APP produce las regiones iniciales, el TRIT fija régimen y
orientación, el TPK produce el registro de transiciones, y ese registro
determina \(L_0,L_1\) y \(0011\) antes de toda realización arquimediana.

